\section{Diagnostica e Configurazione di Rete}

\begin{frame}[fragile]{Verifica della raggiungibilità: \texttt{ping}}
  Il comando \texttt{ping} usa il protocollo ICMP per verificare se un host remoto è attivo e raggiungibile, misurando il tempo di andata e ritorno (\textit{Round Trip Time}, RTT).

  Esecuzione continua (interrompere con \texttt{Ctrl+C}):
  \begin{figure}
    \begin{lstlisting}[backgroundcolor=\color{expowhite!50!white}]
[mattia@pirandello ~]$ ping google.com
\end{lstlisting}
  \end{figure}

  Limitare il numero di pacchetti inviati con il flag \texttt{-c}:
  \begin{figure}
    \begin{lstlisting}[backgroundcolor=\color{expowhite!50!white}]
[mattia@pirandello ~]$ ping -c 4 8.8.8.8
\end{lstlisting}
    \caption{Invia esattamente 4 pacchetti prima di terminare}
  \end{figure}\medskip

  Mostra le statistiche sui pacchetti trasmessi/ricevuti e l'eventuale perdita di rete (\textit{packet loss}).
\end{frame}

\begin{frame}[fragile]{\texttt{ping}: Errori comuni restituiti}
  1. Problema di risoluzione nomi (DNS non configurato o non raggiungibile):
  \begin{figure}
    \begin{lstlisting}[backgroundcolor=\color{expowhite!50!white}]
[mattia@pirandello ~]$ ping hostnonvalido.lan
ping: hostnonvalido.lan: Name or service not known
\end{lstlisting}
  \end{figure}

  2. Host spento, IP errato o pacchetto scartato silenziosamente da un firewall:
  \begin{figure}
    \begin{lstlisting}[backgroundcolor=\color{expowhite!50!white}]
[mattia@pirandello ~]$ ping -c 2 192.168.1.250
PING 192.168.1.250 (192.168.1.250) 56(84) bytes of data.
\end{lstlisting}
  \end{figure}

  3. Manca la rotta locale oppure il router blocca attivamente il transito:
  \begin{figure}
    \begin{lstlisting}[backgroundcolor=\color{expowhite!50!white}]
From 192.168.1.1: icmp_seq=1 Destination Host Unreachable
\end{lstlisting}
  \end{figure}
\end{frame}

\begin{frame}[fragile]{Download di file: \texttt{wget}}
  \texttt{wget} è utile per scaricare file via rete (supporta HTTP, HTTPS e FTP).

  Scaricare un file salvandolo con il suo nome originale:
  \begin{figure}
    \begin{lstlisting}[backgroundcolor=\color{expowhite!50!white}]
[mattia@pirandello ~]$ wget https://example.com/archivio.tar.gz
\end{lstlisting}
  \end{figure}

  Specificare un nome file o percorso di output (\texttt{-O}):
  \begin{figure}
    \begin{lstlisting}[backgroundcolor=\color{expowhite!50!white}]
[mattia@pirandello ~]$ wget -O setup.sh https://example.com/script.sh
\end{lstlisting}
  \end{figure}

  Continuare un download interrotto con il flag \texttt{-c}:
  \begin{figure}
    \begin{lstlisting}[backgroundcolor=\color{expowhite!50!white}]
[mattia@pirandello ~]$ wget -c https://example.com/grande_file.iso
\end{lstlisting}
  \end{figure}
\end{frame}

\begin{frame}[fragile]{Richieste HTTP con \texttt{curl}}
  \texttt{curl} è un tool versatile per trasferire dati da e verso un server. 

  Visualizzare l'output a schermo:
  \begin{figure}
    \begin{lstlisting}[backgroundcolor=\color{expowhite!50!white}]
[mattia@pirandello ~]$ curl https://example.com
\end{lstlisting}
  \end{figure}

  Salvare la risposta su un file:
  \begin{figure}
    \begin{lstlisting}[backgroundcolor=\color{expowhite!50!white}]
[mattia@pirandello ~]$ curl -o pagina.html https://example.com
\end{lstlisting}
  \end{figure}\medskip
\end{frame}

\begin{frame}[fragile]{Richieste HTTP con \texttt{curl}}
  \texttt{curl} è solitamente utilizzato per interagire con le API REST e testare web service.

  Ispezionare solo gli header HTTP di risposta (\texttt{-I}):
  \begin{figure}
    \begin{lstlisting}[backgroundcolor=\color{expowhite!50!white}]
[mattia@pirandello ~]$ curl -I https://example.com
\end{lstlisting}
    \caption{Mostra status code (es. 200 OK), content-type, cookie, ecc.}
  \end{figure}

  Richieste in formato JSON:
  \begin{figure}
    \begin{lstlisting}[backgroundcolor=\color{expowhite!50!white}]
[mattia@pirandello ~]$ curl -X POST https://api.example.com/utenti \
    -H "Content-Type: application/json" \
    -d '{"nome": "Mario", "ruolo": "studente"}'
\end{lstlisting}
  \end{figure}
\end{frame}